\section{Discussion}

The central interpretation is about evaluation units. A random split
can be internally consistent and still answer a less demanding question
than transfer to a new cancer/study cohort. In this corrected six-cohort
matrix, both random and true-patient-grouped CV yield higher
discrimination and lower calibration error than cohort holdout.
The common accession panel was defined with knowledge of which features
occur in the held-out cohort. This is not label leakage, but the
contrast conditions on target-cohort feature availability; an unseen
cohort with a prospectively locked panel remains unevaluated here.
Patient grouping closes the identified patient-identity defect without
erasing the contrast. It does not separate cancer-type biology from
study-specific acquisition conditions.

The distinction between macro and pooled performance is consequential.
Pooled AUROC combines scores from different fitted models and
cross-cohort positive--negative pairs; the macro result averages
within-cohort ranking. Neither should be silently substituted for the
other. Fixed-prediction patient resampling gives a useful sensitivity
to the sampled patients within these six studies, but not a complete
uncertainty analysis of refitted learning procedures or future cohorts.

Global abundance summaries and missingness controls constrain
interpretation of the high-dimensional model without converting AUROC
increments into biological shares. Correct accession identity is
necessary for this comparison; choosing the first protein associated
with each gene does not guarantee a shared protein feature across
cohorts. The source-normalized measurement scale also prevents the
three-summary model from being described as a raw total-signal assay.

The imputation and model-family sensitivities show why similar pooled
performance is insufficient evidence of invariance. L2 and XGBoost
have almost the same pooled AUROC, but differ substantially on CCRCC
and PDAC. Calibration and selective prediction likewise require
cohort-resolved reporting: improved retained accuracy is accompanied
by a changed cohort mixture, and the specified recalibration maps help
some cohorts more than others. Neither supplies a deployment guarantee.

The correction and evaluation questions should remain separate.
BERNN uses unlabeled validation and test features during representation
warm-up \citep{PelletierBERNN}; its reported protocol is not an
interchangeable numerical baseline for the transfer fits here. We have
not benchmarked BERNN or pipeML \citep{HurtadoPipeML} head-to-head and
make no superiority claim. A future harmonization comparison should
state whether unlabeled target features are available and report
calibration as well as discrimination. Our rejection curves are not
conformal prediction sets and have no finite-sample coverage guarantee
under cohort shift.

Tumor-versus-normal is a high-effect testbed, not a clinically actionable
biomarker task. Six selected cohorts from one CPTAC processing source
do not establish multi-platform generality. The separate OV summary is
historical context without prediction-level replay in this revision;
it is not independent evidence for or against the six-cohort contrast.
Historical diagnostics on the superseded feature matrix remain
archived but cannot substitute for new evidence after identities change.
Independent measurements crossing cancer type with acquisition
environment, broader model comparisons, and full refit-based uncertainty
would be needed for stronger transfer claims.

The practical recommendation is to align the evaluation unit with the
intended use, report within-cohort discrimination, calibration and
coverage together, and distinguish macro from pooled summaries.
The corrected data still support a protocol contrast in this design.
They do not establish a general biological transfer law or identify a
causal batch effect.

\section{Supporting Information}
The accompanying Supporting Information gives corrected sample and
feature definitions, cohort-resolved AUROC and probability errors,
conditional resampling ranges, L2 convergence checks, and an explicit
disposition of historical side studies. Its tables and the main-text
tables are generated from the same verified prediction records.

\section{Data Availability Statement}
All biological inputs are public CPTAC tables accessed through
\texttt{cptac} \citep{cptacPkg}. The primary analysis uses the
\texttt{umich} source for CCRCC, LUAD, UCEC, HNSCC, LSCC and PDAC.
No new biological data were generated.

The current code and research artifacts are available at
\url{https://github.com/PeterPonyu/b28-proteomics-transfer}, release tag
\texttt{author-review-2026-10-02}, with the version-specific archive
\url{https://doi.org/10.5281/zenodo.23101188}. The package contains this
Article and Supporting Information, their source and figure inputs,
corrected sample/patient/fold and accession maps, stored predictions,
analysis and verification programs, software environment records and a
SHA-256 inventory. Dated candidate filenames retain their provenance;
the current public manifest binds the included files to this revision.
Default offline checks verify stored predictions and metrics and render
the supplied PDFs without fitting models or downloading assay data.

A complete scientific rerun requires separately acquired, matching
CPTAC/UMICH abundance and mapping files and the recorded fit environment.
\texttt{DATA\_SOURCES.md} records cohort/source/loader versions, filenames
and input hashes; unrecorded PDC study accessions and per-cohort version
labels are explicitly identified as unknown. Third-party raw proteomics
are not redistributed. Author-owned code is MIT licensed; authored text,
figures and derived result tables are CC BY 4.0 to the extent of author
rights. Upstream conditions remain applicable as detailed in
\texttt{THIRD\_PARTY\_NOTICES.md}. The earlier archive
\url{https://doi.org/10.5281/zenodo.21870024} is historical, not the
current corrected release. Artifact publication does not constitute
independent scientific or clinical validation.
